\sloppy
\chapter{Related Work}
\addtocontents{toc}{\protect\setcounter{tocdepth}{2}}
\section{Introduction}

\citep{Kaliyar2021FakeBERT} proposed FakeBERT, a BERT-based model for fake news detection on social media. It adds parallel CNN layers to the BERT architecture for multi-scale feature extraction and was tested on the LIAR and ISOT datasets. The model reached 98.90\% accuracy, captured contextual semantic relationships well, and clearly outperformed traditional ML models. It is, however, expensive to run, needs a large amount of GPU memory, and handles only English. The hybrid model in this thesis follows the same transformer-centred idea but places a BiLSTM on top of BERT's contextual embeddings instead of parallel CNNs, and it is evaluated on several datasets rather than one.

\citep{Aphiwongsophon2021Detecting} compared several classical machine learning classifiers for fake news detection. They used linguistic and statistical features extracted with TF-IDF and trained Logistic Regression, SVM, Naive Bayes, and Random Forest. Such models are simple to implement, cheap to run, and interpretable, which suits resource-constrained settings. Their accuracy trails deep learning because they cannot capture semantic context or word order. The system proposed here relies on BERT's contextual embeddings instead, giving up some of that simplicity for a much richer representation of meaning.

\citep{Nasir2021FakeNews} proposed a hybrid deep learning architecture that combines a CNN for local feature extraction with an RNN for sequential pattern learning. It was trained with GloVe word embeddings on the ISOT and FA-KES datasets. The model captures local and sequential features together and generalizes better than single-model approaches, but it needs a lot of training data, infers slowly, and struggles with short text such as tweets and headlines. Our BERT + BiLSTM model addresses the short-text weakness by using BERT's pre-trained contextual representations instead of embeddings learned from scratch.

\citep{Przybyla2020Capturing} analysed stylometric and linguistic style features that separate fake from credible news, such as sentence length, punctuation patterns, and readability scores, using an SVM classifier. The approach needs no external knowledge base, works on writing style alone, and is light and fast at inference. Its main weakness is that a skilled fake news creator can easily imitate the style, so it is not robust against adversarial content. Contextual, semantics-driven models such as the one proposed here are comparatively less exposed to this problem.

\citep{Popat2019DeClarE} proposed an evidence-aware approach that retrieves external web evidence and combines it with claim credibility assessment through an attention-based LSTM, evaluated on the Snopes and PolitiFact datasets. It brings in external knowledge for verification, its attention mechanism gives interpretability, and it performs well on political claims. It depends on evidence being available, is slow because of web retrieval, and fails when evidence is missing or biased. The present work does no evidence retrieval and concentrates on the classification model itself, although Chapter~5 names retrieval-augmented evidence grounding as a natural extension.

\citep{Shu2020FakeNewsNet} built FakeNewsNet, a fake news benchmark that includes news content, user social engagement, and spatiotemporal propagation features from PolitiFact and GossipCop. It is among the most comprehensive public datasets and supports multi-modal detection research. Collecting it is time-consuming, and its social-context features need platform API access, which is increasingly restricted. This thesis uses the more text-centric LIAR and ISOT datasets, which suit a controlled, text-only model comparison better.

\citep{Mehta2022RecentTrends} surveyed state-of-the-art deep learning techniques for fake news detection, covering transformer models, graph neural networks, and multi-modal approaches. The survey covers recent methods broadly and points out research gaps, but it proposes no new model and includes only English-language studies. The systematic comparison of sequential, transformer, and hybrid models in this thesis addresses one of the concrete gaps that a survey like this highlights.

\citep{Rai2022FakeNews} combined BERT contextual embeddings with an enhanced LSTM and an additional attention mechanism, and trained the model on the LIAR and ISOT datasets. The hybrid outperformed standalone BERT and standalone LSTM, and the attention layer improved interpretability and robustness to variation in text length. The drawbacks are high model complexity, longer training, and the need for a lot of labelled data for fine-tuning. This is the closest published work to the BERT + BiLSTM model proposed here. This thesis leaves out the extra attention layer so that it can isolate what sequential modelling adds on top of BERT.

\citep{CarrascoFarre2022Fingerprints} studied psychological and cognitive bias patterns in fake news content, using computational linguistics and psycholinguistic feature extraction with the LIWC tool together with ML classifiers. This gives a new view that joins psychology and NLP, with interpretable features that apply across domains and languages. The study is limited by a small dataset and by psycholinguistic tools that may not carry over to other cultural and linguistic contexts. This line of analysis complements the purely text-based transformer models compared here and lies outside the scope of this thesis.

\citep{Huang2020FakeNews} proposed an optimized ensemble model that uses a metaheuristic harmony search algorithm for automatic feature selection and classifier weight optimization. It combines SVM, Random Forest, and Logistic Regression on several news datasets. Automatic feature selection reduces dimensionality, and the ensemble is more robust and outperforms the individual classifiers. The metaheuristic step adds heavy computation, and the implementation is too complex for real-time detection. The transformer-based models evaluated in this thesis avoid this optimization overhead by relying on pre-trained representations.

\citep{Shu2017FakeNewsDetection} gave a foundational data mining perspective on fake news detection on social media. It proposed a taxonomy of detection methods covering knowledge-based, style-based, propagation-based, and credibility-based approaches. The taxonomy spans many detection paradigms and remains a standard reference. As a survey it has no experimental results and covers deep learning and transformer-based models only to a limited extent. This thesis falls in the credibility-based, content-driven branch of that taxonomy.

\citep{Zhou2020Survey} surveyed the fundamental theories, detection methods, and future opportunities in fake news research. They reviewed knowledge-based, style-based, propagation-based, and credibility-based methods across several benchmark datasets. The survey covers detection strategies thoroughly and identifies open problems, but it proposes no detection model and gives little space to multimodal and LLM-based approaches. These gaps motivate the future work in Chapter~5 of this thesis.

\citep{Ruchansky2017CSI} proposed CSI (Capture, Score, Integrate), a hybrid deep learning model that combines text content features with user behaviour and propagation patterns. It uses a hybrid RNN architecture and was evaluated on a CIKM benchmark dataset. CSI was the first model to integrate content, user, and propagation signals, and it clearly outperformed text-only baselines on social media platforms. It needs social engagement data that may not exist, has a complex multi-module architecture, and generalizes poorly to news outside social media. The present work stays text-only on purpose, giving up the extra signal that CSI uses in return for a simpler and more widely applicable model.

\citep{Vosoughi2018Spread} carried out a landmark empirical study of how true and false news spread on Twitter. They analysed 126,000 news stories tweeted by three million users over ten years and found that false news spreads about six times faster than true news and reaches more people. It is the largest study of its kind and strongly motivates automated detection. It is purely observational, proposes no detection model, and uses only Twitter data. This finding is one of the central motivations cited in Chapter~1.

\citep{Wang2017Liar} introduced the LIAR benchmark, which contains 12,836 manually labelled short statements from PolitiFact with six fine-grained truthfulness labels. It comes with metadata such as speaker, subject, context, and party affiliation, and with baseline CNN and BiLSTM classifiers. LIAR was the first large-scale manually labelled fake news dataset, and its rich metadata supports multi-feature modelling. It is limited to political speech, and the fine-grained labels make classification harder. LIAR is one of the two benchmark datasets used to evaluate the models in this thesis.

\citep{Touvron2023Llama2} introduced LLaMA 2, a family of open-source large language models from 7B to 70B parameters, fine-tuned for dialogue and instruction following and evaluated against GPT-4 and other proprietary models. It is free for research, competitive with closed models, and allows privacy-preserving local deployment. It needs substantial GPU resources, can still hallucinate, and has weaker multilingual ability than GPT-4. Chapter~5 names fine-tuning a model such as LLaMA 2 on fake news data as a future direction.

\citep{Lewis2020RAG} proposed Retrieval-Augmented Generation (RAG), a framework that combines the parametric knowledge of pre-trained LLMs with non-parametric retrieval from external knowledge bases. It pairs dense passage retrieval with a sequence-to-sequence generator and was evaluated on open-domain QA, fact verification, and knowledge-grounded dialogue. RAG lets an LLM use up-to-date external evidence and reduces hallucination without retraining the whole model. Its retrieval quality depends heavily on how well the knowledge base covers the topic, and it adds inference latency. Chapter~5 proposes grounding fake news classification decisions in evidence retrieved by a RAG-style pipeline as future work.

\citep{Ribeiro2016LIME} proposed LIME (Local Interpretable Model-agnostic Explanations), which explains the prediction of any black-box classifier by approximating it locally with an interpretable model. It works with any classifier regardless of architecture and produces human-readable explanations that support trust and debugging. Its local explanations may not reflect the model's global behaviour, and it can be expensive for large input spaces. Chapter~5 discusses applying LIME to show which words drove a BERT-based REAL or FAKE prediction.

\citep{Lundberg2017SHAP} proposed SHAP (SHapley Additive exPlanations), a unified framework that interprets model predictions with Shapley values from cooperative game theory. It is theoretically grounded, consistent, and locally accurate, and it supports both global feature importance and explanations of single predictions. It is expensive for large feature spaces and deep models and needs careful interpretation. Chapter~5 also names the combination of SHAP with attention maps, to quantify word-level contributions in the BERT-based classifier, as future work.

\section{Comparative Summary of the Literature}

Table~\ref{tab:litreview} summarizes the nineteen studies above by contribution, method, advantage, disadvantage, and the future scope each one names. It shows the pattern that motivates this thesis: classical and stylometric models are cheap but shallow, transformer and hybrid models are accurate but costly, and few of the surveyed studies test their models on a dataset other than the one they were trained on.

{\footnotesize
\begin{longtable}{|>{\raggedright\arraybackslash}p{0.17\textwidth}|>{\raggedright\arraybackslash}p{0.17\textwidth}|>{\raggedright\arraybackslash}p{0.17\textwidth}|>{\raggedright\arraybackslash}p{0.17\textwidth}|>{\raggedright\arraybackslash}p{0.17\textwidth}|}
\caption{Comparative Literature Review of Fake News Detection Methods} \label{tab:litreview} \\
\hline
\textbf{Contribution} & \textbf{Method} & \textbf{Advantage} & \textbf{Disadvantage} & \textbf{Future Scope} \\ \hline
\endfirsthead
\hline
\textbf{Contribution} & \textbf{Method} & \textbf{Advantage} & \textbf{Disadvantage} & \textbf{Future Scope} \\ \hline
\endhead
\hline
\multicolumn{5}{r}{\small\textit{Continued on next page\ldots}} \\
\endfoot
\hline
\endlastfoot
Introduced FakeBERT, a BERT-based deep learning model specifically designed for fake news detection on social media platforms with parallel CNN layers integrated into BERT architecture \citep{Kaliyar2021FakeBERT}. &
BERT + Parallel CNN layers for multi-scale feature extraction on LIAR and ISOT datasets \citep{Kaliyar2021FakeBERT}. &
Achieved 98.90\% accuracy; effectively captures contextual semantic relationships in text; outperforms traditional ML models significantly \citep{Kaliyar2021FakeBERT}. &
Computationally expensive; requires large GPU memory; limited to English language content only \citep{Kaliyar2021FakeBERT}. &
Extend to multilingual detection; reduce computational cost using model distillation techniques \citep{Kaliyar2021FakeBERT}. \\ \hline
Systematically compared multiple classical ML classifiers for fake news detection using linguistic and statistical text features extracted via TF-IDF \citep{Aphiwongsophon2021Detecting}. &
Logistic Regression, SVM, Naive Bayes, Random Forest with TF-IDF feature vectors on benchmark news datasets \citep{Aphiwongsophon2021Detecting}. &
Simple implementation; low computational cost; interpretable results; suitable for resource-constrained environments \citep{Aphiwongsophon2021Detecting}. &
Lower accuracy compared to deep learning; cannot capture semantic context or word order in text \citep{Aphiwongsophon2021Detecting}. &
Combine classical ML with deep learning features; incorporate metadata and social context signals \citep{Aphiwongsophon2021Detecting}. \\ \hline
Proposed a hybrid deep learning architecture combining CNN for local feature extraction and RNN for sequential pattern learning in news text \citep{Nasir2021FakeNews}. &
Hybrid CNN-RNN architecture with GloVe word embeddings trained on ISOT and FA-KES datasets \citep{Nasir2021FakeNews}. &
Captures both local and sequential text features simultaneously; better generalization than single-model approaches \citep{Nasir2021FakeNews}. &
Requires significant training data; slow inference time; struggles with short text like tweets and headlines \citep{Nasir2021FakeNews}. &
Apply to real-time social media streams; incorporate attention mechanisms for better interpretability \citep{Nasir2021FakeNews}. \\ \hline
Analyzed stylometric and linguistic style features that distinguish fake news from credible news using writing patterns and rhetorical structures \citep{Przybyla2020Capturing}. &
Stylometric analysis with SVM classifier; features include sentence length, punctuation patterns, readability scores on news datasets \citep{Przybyla2020Capturing}. &
Does not require external knowledge base; works purely on writing style; lightweight and fast inference \citep{Przybyla2020Capturing}. &
Style-based features are easily mimicked by sophisticated fake news creators; low robustness against adversarial content \citep{Przybyla2020Capturing}. &
Combine style features with semantic embeddings for more robust multi-feature detection systems \citep{Przybyla2020Capturing}. \\ \hline
Proposed evidence-aware fake news detection by retrieving external web evidence and combining it with claim credibility assessment using attention networks \citep{Popat2019DeClarE}. &
Attention-based LSTM with external web evidence retrieval and source credibility scoring on Snopes and PolitiFact datasets \citep{Popat2019DeClarE}. &
Incorporate external knowledge for verification; attention mechanism provides interpretability; better performance on political claims \citep{Popat2019DeClarE}. &
Dependent on external evidence availability; slow due to web retrieval; fails when evidence is unavailable or biased \citep{Popat2019DeClarE}. &
Integrate real-time knowledge graphs; improve evidence retrieval with more reliable fact-checking APIs \citep{Popat2019DeClarE}. \\ \hline
Developed a comprehensive fake news benchmark dataset incorporating news content, user social engagement, and spatiotemporal propagation features \citep{Shu2020FakeNewsNet}. &
Multi-modal dataset construction combining news text, social network graphs, user profiles, and publisher credibility on PolitiFact and GossipCop \citep{Shu2020FakeNewsNet}. &
Most comprehensive publicly available dataset; enables multi-modal detection research; includes social context signals \citep{Shu2020FakeNewsNet}. &
Dataset collection is time-consuming; social context features require platform API access which is increasingly restricted \citep{Shu2020FakeNewsNet}. &
Expand dataset to include video and image modalities; add multilingual fake news samples \citep{Shu2020FakeNewsNet}. \\ \hline
Surveyed state-of-the-art deep learning techniques for fake news detection covering transformer models, graph neural networks, and multi-modal approaches \citep{Mehta2022RecentTrends}. &
Systematic review of BERT, GPT, GNN, and multi-modal models across multiple benchmark datasets from 2017 to 2022 \citep{Mehta2022RecentTrends}. &
Comprehensive coverage of recent methods; identifies research gaps; provides clear taxonomy of detection approaches \citep{Mehta2022RecentTrends}. &
Survey only; no novel model proposed; coverage limited to English language studies \citep{Mehta2022RecentTrends}. &
Develop unified evaluation benchmark; explore cross-lingual and low-resource fake news detection \citep{Mehta2022RecentTrends}. \\ \hline
Combined BERT contextual embeddings with enhanced LSTM architecture for improved sequential and contextual understanding of fake news text \citep{Rai2022FakeNews}. &
BERT + enhanced LSTM hybrid model with attention mechanism trained on LIAR and ISOT datasets \citep{Rai2022FakeNews}. &
Superior performance over standalone BERT or LSTM; attention layer improves interpretability; robust to text length variations \citep{Rai2022FakeNews}. &
High model complexity; longer training time; requires substantial labeled data for effective fine-tuning \citep{Rai2022FakeNews}. &
Apply transfer learning to low-resource languages; integrate image features for multi-modal detection \citep{Rai2022FakeNews}. \\ \hline
Investigated psychological and cognitive bias patterns embedded in fake news content using computational linguistics and psycholinguistic feature analysis \citep{CarrascoFarre2022Fingerprints}. &
Psycholinguistic feature extraction using LIWC tool with ML classifiers; analysis of cognitive bias markers in fake vs.\ real news articles \citep{CarrascoFarre2022Fingerprints}. &
Novel perspective combining psychology and NLP; interpretable features; applicable across domains and languages \citep{CarrascoFarre2022Fingerprints}. &
Limited dataset size; psycholinguistic tools may not generalize across cultural and linguistic contexts \citep{CarrascoFarre2022Fingerprints}. &
Extend psycholinguistic analysis to social media posts; develop cross-cultural misinformation detection models \citep{CarrascoFarre2022Fingerprints}. \\ \hline
Proposed an optimized ensemble learning model using metaheuristic harmony search algorithm for automatic feature selection and classifier weight optimization \citep{Huang2020FakeNews}. &
Ensemble of SVM, RF, and LR optimized using Self-Adaptive Harmony Search metaheuristic algorithm on multiple news datasets \citep{Huang2020FakeNews}. &
Automatic feature selection reduces dimensionality; ensemble approach improves robustness; outperforms individual classifiers \citep{Huang2020FakeNews}. &
High computational overhead due to metaheuristic optimization; complex implementation; not suitable for real-time detection \citep{Huang2020FakeNews}. &
Apply neural architecture search for automated deep learning model optimization in fake news detection \citep{Huang2020FakeNews}. \\ \hline
Provided a foundational data mining perspective on fake news detection on social media, proposing taxonomy of detection methods including knowledge-based, style-based, propagation-based, and credibility-based approaches \citep{Shu2017FakeNewsDetection}. &
Survey and taxonomy of fake news detection approaches on social media platforms; analysis of linguistic, visual, and social context features \citep{Shu2017FakeNewsDetection}. &
Comprehensive taxonomy; covers diverse detection paradigms; foundational reference for the research community \citep{Shu2017FakeNewsDetection}. &
Survey only; no experimental results; limited coverage of deep learning and transformer-based models \citep{Shu2017FakeNewsDetection}. &
Develop end-to-end multi-modal systems combining all four detection paradigms for improved accuracy \citep{Shu2017FakeNewsDetection}. \\ \hline
Presented a comprehensive survey of fake news fundamental theories, detection methods, and future opportunities covering the full spectrum from fabrication to verification \citep{Zhou2020Survey}. &
Systematic review of knowledge-based, style-based, propagation-based, and credibility-based detection methods across multiple benchmark datasets \citep{Zhou2020Survey}. &
Thorough coverage of detection strategies; identifies open problems; provides structured framework for future research \citep{Zhou2020Survey}. &
Survey only; does not propose a novel detection model; limited coverage of multimodal and LLM-based approaches \citep{Zhou2020Survey}. &
Extend survey coverage to include graph neural networks, multimodal detection, and cross-lingual fake news models \citep{Zhou2020Survey}. \\ \hline
Proposed CSI (Capture, Score, Integrate), a hybrid deep learning model combining text content features with user behaviour and propagation patterns for fake news detection \citep{Ruchansky2017CSI}. &
Hybrid RNN model integrating article text (Capture), user response scoring, and propagation network modelling on CIKM benchmark dataset \citep{Ruchansky2017CSI}. &
First to integrate content, user, and propagation signals; significantly outperforms text-only baselines; applicable to social media platforms \citep{Ruchansky2017CSI}. &
Requires social engagement data which may be unavailable; complex multi-module architecture; limited generalization to non-social media news \citep{Ruchansky2017CSI}. &
Extend to real-time streaming data; integrate image and video modalities; explore graph-based propagation modelling \citep{Ruchansky2017CSI}. \\ \hline
Conducted a landmark empirical study on the spread of true and false news online across Twitter, revealing that false news spreads six times faster than true news and reaches more people \citep{Vosoughi2018Spread}. &
Large-scale analysis of 126,000 news stories tweeted by three million users over ten years; statistical modelling of diffusion cascades \citep{Vosoughi2018Spread}. &
Largest empirical study of its kind; provides strong motivation for automated fake news detection; highlights human-driven amplification of misinformation \citep{Vosoughi2018Spread}. &
Observational study; does not propose a detection model; limited to Twitter platform data \citep{Vosoughi2018Spread}. &
Extend analysis to other social media platforms; design detection systems that model propagation speed as a real-time detection signal \citep{Vosoughi2018Spread}. \\ \hline
Introduced the LIAR benchmark dataset containing 12,836 manually labelled short statements from PolitiFact with six fine-grained truthfulness labels for fake news detection research \citep{Wang2017Liar}. &
Benchmark dataset construction with metadata including speaker, subject, context, and party affiliation; baseline CNN and BiLSTM classifiers evaluated \citep{Wang2017Liar}. &
First large-scale manually labelled fake news dataset; enables fine-grained classification; metadata-rich annotations support multi-feature modelling \citep{Wang2017Liar}. &
Limited to political speech domain; fine-grained labels make classification harder; dataset may not generalize to social media news content \citep{Wang2017Liar}. &
Extend dataset with image and video evidence; develop cross-domain transfer learning models trained on LIAR for general fake news detection \citep{Wang2017Liar}. \\ \hline
Introduced LLaMA 2, a family of open-source large language models ranging from 7B to 70B parameters, enabling state-of-the-art language understanding and generation for downstream NLP tasks \citep{Touvron2023Llama2}. &
Open-source transformer-based LLM family fine-tuned for dialogue and instruction following; evaluated on multiple NLP benchmarks against GPT-4 and other proprietary models \citep{Touvron2023Llama2}. &
Freely available for research; competitive performance against closed models; enables privacy-preserving local deployment of large language models \citep{Touvron2023Llama2}. &
Requires significant GPU resources for inference and fine-tuning; may still produce hallucinations; limited multilingual capability compared to GPT-4 \citep{Touvron2023Llama2}. &
Fine-tune LLaMA 2 on fake news datasets; integrate with RAG pipelines to enable evidence-grounded fake news classification \citep{Touvron2023Llama2}. \\ \hline
Proposed Retrieval-Augmented Generation (RAG), a framework combining parametric knowledge of pre-trained LLMs with non-parametric retrieval from external knowledge bases for knowledge-intensive NLP tasks \citep{Lewis2020RAG}. &
RAG model combining dense passage retrieval with a sequence-to-sequence generator; evaluated on open-domain QA, fact verification, and knowledge-grounded dialogue tasks \citep{Lewis2020RAG}. &
Enables LLMs to access up-to-date external evidence; reduces hallucination; improves factual accuracy without full model retraining \citep{Lewis2020RAG}. &
Retrieval quality is highly dependent on the quality and coverage of the knowledge base; increased inference latency due to retrieval step \citep{Lewis2020RAG}. &
Apply RAG to fake news detection by grounding classification decisions in retrieved evidence from verified fact-checking databases \citep{Lewis2020RAG}. \\ \hline
Proposed LIME (Local Interpretable Model-agnostic Explanations), a technique for explaining the predictions of any black-box classifier by approximating it locally with an interpretable model \citep{Ribeiro2016LIME}. &
Model-agnostic local surrogate modelling using perturbed samples to identify which input features most influenced a specific prediction across text, image, and tabular data \citep{Ribeiro2016LIME}. &
Works with any classifier regardless of architecture; provides human-readable explanations; enables trust and debugging of black-box ML models \citep{Ribeiro2016LIME}. &
Local explanations may not reflect global model behaviour; computationally expensive for large input spaces; sensitive to perturbation sampling strategy \citep{Ribeiro2016LIME}. &
Apply LIME to explain fake news classification decisions; integrate with BERT-based models to highlight which words drove the REAL or FAKE prediction \citep{Ribeiro2016LIME}. \\ \hline
Proposed SHAP (SHapley Additive exPlanations), a unified framework for interpreting model predictions using Shapley values from cooperative game theory to quantify each feature's contribution \citep{Lundberg2017SHAP}. &
Unified additive feature attribution method grounded in Shapley values; applied to tree-based models, deep learning, and linear models for consistent global and local explanations \citep{Lundberg2017SHAP}. &
Theoretically grounded in game theory; consistent and locally accurate; enables both global feature importance and individual prediction explanations \citep{Lundberg2017SHAP}. &
Computationally expensive for large feature spaces and deep models; requires careful interpretation; may not capture feature interaction effects completely \citep{Lundberg2017SHAP}. &
Apply SHAP to quantify word-level contributions in BERT fake news classifiers; combine with attention maps for multi-level explainability dashboards \citep{Lundberg2017SHAP}. \\ \hline
\end{longtable}
}

\section{Summary}

The literature reviewed here covers classical machine learning classifiers, hybrid deep learning architectures, evidence-aware and psycholinguistic approaches, benchmark dataset papers, broad surveys, and recent work on large language models and explainability. One pattern runs through it. Classical statistical methods stay cheap and interpretable but cannot capture semantic context, while deep learning and transformer-based methods reach higher accuracy at a greater computational cost. Few of the reviewed works compare sequential, transformer-based, and hybrid architectures under the same controlled conditions. We believe a systematic comparison of a standalone BiLSTM, a standalone BERT model, and a hybrid BERT + BiLSTM model, evaluated consistently on several benchmark datasets, addresses this gap. It also gives clearer practical guidance on when the added complexity of a hybrid architecture is justified.
